\documentclass[10pt, letterpaper]{article}

% Packages:
\usepackage[
    ignoreheadfoot,
    top=0.1 cm,
    bottom=0.1 cm,
    left=0.1 cm,
    right=0.1 cm
]{geometry}

\usepackage{titlesec}
\usepackage{tabularx}
\usepackage{array}
\usepackage[dvipsnames]{xcolor}
\definecolor{primaryColor}{RGB}{0, 79, 144}
\usepackage{enumitem}
\usepackage{fontawesome5}
\usepackage{amsmath}
\usepackage[
    pdftitle={Ronit Gandhi's CV},
    pdfauthor={Ronit Gandhi},
    pdfcreator={LaTeX with RenderCV},
    colorlinks=true,
    urlcolor=primaryColor
]{hyperref}
\usepackage{setspace}
\setstretch{0.7}
\usepackage[pscoord]{eso-pic}
\usepackage{calc}
\usepackage{bookmark}
\usepackage{lastpage}
\usepackage{changepage}
\usepackage{paracol}
\usepackage{ifthen}
\usepackage{needspace}
\usepackage{iftex}

% Ensuring machine readability
\ifPDFTeX
    \input{glyphtounicode}
    \pdfgentounicode=1
    \usepackage[utf8]{inputenc}
    \usepackage{lmodern}
\fi

% Style adjustments
\AtBeginEnvironment{adjustwidth}{\partopsep0pt}
\pagestyle{empty}
\setcounter{secnumdepth}{0}
\setlength{\parindent}{0pt}
\setlength{\topskip}{0pt}
\setlength{\columnsep}{0cm}
\makeatletter
\let\ps@customFooterStyle\ps@plain
\patchcmd{\ps@customFooterStyle}{\thepage}{}{}{}
\makeatother
\pagestyle{customFooterStyle}

\titleformat{\section}{\needspace{4\baselineskip}\bfseries\large}{}{0pt}{}[\vspace{1pt}\titlerule]
\titlespacing{\section}{-1pt}{0.3 cm}{0.2 cm}

\renewcommand\labelitemi{$\circ$}

\newenvironment{highlights}{
    \begin{itemize}[
        topsep=0.10 cm,
        parsep=0.10 cm,b
        partopsep=0pt,
        itemsep=0pt,
        leftmargin=0.4 cm + 10pt
    ]
}{
    \end{itemize}
}

\newenvironment{highlightsforbulletentries}{
    \begin{itemize}[
        topsep=0.10 cm,
        parsep=0.10 cm,
        partopsep=0pt,
        itemsep=0pt,
        leftmargin=10pt
    ]
}{
    \end{itemize}
}

\newenvironment{onecolentry}{
    \begin{adjustwidth}{
        0.2 cm + 0.00001 cm
    }{
        0.2 cm + 0.00001 cm
    }
}{
    \end{adjustwidth}
}

\newenvironment{twocolentry}[2][]{
    \onecolentry
    \def\secondColumn{#2}
    \setcolumnwidth{\fill, 4.5 cm}
    \begin{paracol}{2}
}{
    \switchcolumn \raggedleft \secondColumn
    \end{paracol}
    \endonecolentry
}

\newenvironment{header}{
    \setlength{\topsep}{0pt}\par\kern\topsep\centering\linespread{1.5}
}{
    \par\kern\topsep
}

\newcommand{\placelastupdatedtext}{%
  \AddToShipoutPictureFG*{
    \put(\LenToUnit{\paperwidth-2 cm-0.2 cm+0.05cm},\LenToUnit{\paperheight-1.0 cm}){\vtop{{\null}\makebox[0pt][c]{\small\color{gray}\textit{}}}}
  }
}

\let\hrefWithoutArrow\href
\renewcommand{\href}[2]{\hrefWithoutArrow{#1}{\ifthenelse{\equal{#2}{}}{ }{#2 }\raisebox{.15ex}{\footnotesize \faExternalLink*}}}

\begin{document}

\placelastupdatedtext

\begin{header}
    \textbf{\fontsize{24 pt}{24 pt}\selectfont Ronit Gandhi}

    \vspace{0.3 cm}

    \normalsize
    \mbox{{\color{black}\footnotesize\faMapMarker*}\hspace*{0.12cm}New Jersey}%
    \kern 0.25 cm%
    \mbox{\hrefWithoutArrow{mailto:ronit.gandhi@rutgers.edu}{\color{black}{\footnotesize\faEnvelope[regular]}\hspace*{0.12cm}ronit.gandhi@rutgers.edu}}%
    \kern 0.25 cm%
    \mbox{\hrefWithoutArrow{tel:+1-732-470-6368}{\color{black}{\footnotesize\faPhone*}\hspace*{0.13cm}+1 732 470 6368}}%
    \kern 0.25 cm%
    \mbox{\hrefWithoutArrow{https://www.linkedin.com/in/ronit-gandhi-5a65941aa/}{\color{black}{\footnotesize\faLinkedinIn}\hspace*{0.12cm}www.linkedin.com/in/ronit-gandhi-5a65941aa}}%
    \kern 0.25 cm%
    \mbox{\hrefWithoutArrow{https://github.com/RonitGandhi}{\color{black}{\footnotesize\faGithub}\hspace*{0.12cm}RonitGandhi}}%
\end{header}


\section{Education}

\begin{twocolentry}{\textit{May 2025}}
    \textbf{Rutgers University} \\
    \textit{M.S. in Statistics \& Data Science} \\
    {\footnotesize \textit{Relevant Coursework: Machine Learning, NLP, Statistical Modeling, Database Management}}
\end{twocolentry}

\begin{twocolentry}{\textit{May 2023}}
    \textbf{University of Mumbai} \\
    \textit{B.Tech. in Electronics Engineering, Minor in Computer Science}
\end{twocolentry}

% \section{Work Experience}

% \begin{twocolentry}{
%     \textit{New Brunswick, NJ}\\
%     \textit{Jun 2024 – Present}
% }
%     \textbf{Machine Learning Engineer} \\
%     \textit{Particulate Systems Laboratory, Rutgers School of Engineering}
% \end{twocolentry}

% \begin{onecolentry}

% \begin{highlights}

%   \item Built and validated statistical \& ML models (Random Forests, GLMs, LSTM, Neural Networks) in Python \& MATLAB for biomass forecasting (R²\textgreater 0.92), with calibrated prediction intervals.  
%   \item Integrated probabilistic forecasts from complex datasets into a real-time MPC framework—using uncertainty quantification to automate setpoints and boost yield efficiency by 15\%. 
  
%   \item Applied analytical thinking through transfer learning and Bayesian hyperparameter optimization (Optuna) to enhance model robustness under variability and sensor noise.
%   \item Engineered hybrid PBE–ML surrogates for granulation control, achieving 4× faster simulations while preserving key metrics (d\textsubscript{43}, d\textsubscript{width}, F\textsubscript{dmass}).  
%   \item Took a pragmatic approach for R\&D of NLP pipelines—fuzzy matching \& TF-IDF for conflict-of-interest detection across 10K+ proposals, and fine-tuned BART/T5/Pegasus for summarization (ROUGE-2+12–15\%). 
%   \item Collaborated with cross-functional researchers to co-author two peer-reviewed papers (Control Eng. Practice, Array) and lead an LLM-summarization manuscript, effectively communicating statistical findings (ANOVA, confidence intervals).



% \end{highlights}
% \end{onecolentry}

% \vspace{0.2 cm}

% \begin{twocolentry}{
%     \textit{New Brunswick, NJ}\\
%     \textit{Jun 2024 – Aug 2024}
% }
%     \textbf{Research Assistant - Data Scientist} \\
%     \textit{Under Prof. Jerry Kim, Rutgers Business School}
% \end{twocolentry}

% \begin{onecolentry}
% \begin{highlights}
%     \item Applied machine learning methods using Python, R, and EconML to analyze large baseball performance datasets. Maintained scalable data pipelines and documented reproducible research workflows.
%     \item Effectively communicated statistical insights to stakeholders, demonstrating people skills \& enabling strategic decisions.
% \end{highlights}
% \end{onecolentry}

% \vspace{0.2 cm}
\section{Work Experience}

\begin{twocolentry}{
    \textit{Santa Clara, CA}\\
    \textit{Aug 2024 – Present}
}
    \textbf{Software Engineering Intern} \\
    \textit{Exo Imaging Inc.}
\end{twocolentry}

\begin{onecolentry}
\begin{highlights}
\item Developed scalable REST APIs using \textbf{Python}, \textbf{Django}, and \textbf{PostgreSQL (pgvector)} for intelligent medical imaging search and retrieval with low-latency performance.
\item Integrated and fine-tuned \textbf{LLM endpoints (AWS Bedrock: Titan, Nova Pro, Cohere)} for \textbf{text-to-SQL} and semantic search, improving model accuracy and query responsiveness in production.
\item Deployed and optimized AI workflows on \textbf{AWS (EC2, Lambda, S3, RDS)} with performance tuning and observability tools to ensure seamless, scalable model integration.



\end{highlights}
\end{onecolentry}

\vspace{0.2 cm}

\begin{twocolentry}{
    \textit{New Brunswick, NJ}\\
    \textit{Jun 2024 – July 2025}
}
    \textbf{Machine Learning Engineer} \\
    \textit{Particulate Systems Laboratory, Rutgers School of Engineering}
\end{twocolentry}

\begin{onecolentry}
\begin{highlights}
  \item Built and validated statistical \& ML models (Random Forests, GLMs, LSTM, Neural Networks) in Python \& MATLAB for biomass forecasting (R²\textgreater 0.92).  
  \item Integrated probabilistic forecasts into a real-time MPC framework, using uncertainty quantification to boost yield efficiency by 15\%.  
  \item Applied Bayesian hyperparameter optimization (Optuna) and transfer learning to enhance model robustness under variability and sensor noise.
  \item Developed hybrid PBE–ML surrogates for process control, achieving 4× faster simulations while preserving accuracy of key quality metrics.
\end{highlights}
\end{onecolentry}

\vspace{0.2 cm}

\begin{twocolentry}{
    \textit{New Brunswick, NJ}\\
    \textit{Jun 2024 – Aug 2024}
}
    \textbf{Research Assistant - Data Scientist} \\
    \textit{Under Prof. Jerry Kim, Rutgers Business School}
\end{twocolentry}

\begin{onecolentry}
\begin{highlights}
    \item Applied machine learning methods using Python, R, and EconML to analyze large baseball performance datasets. Maintained scalable data pipelines and documented reproducible research workflows.
    \item Effectively communicated statistical insights to stakeholders, enabling strategic decisions.
\end{highlights}
\end{onecolentry}

\vspace{0.2 cm}

\begin{twocolentry}{
    \textit{Mumbai, India}\\
    \textit{Oct 2021 – Jun 2022}
}
    \textbf{Data Scientist Intern, Google via DKSH Smollan}
\end{twocolentry}

\begin{onecolentry}
\begin{highlights}
    \item Built a Python-based BI tool on GCP (BigQuery, Cloud Functions) with LookerStudio and Selenium scraping, improving data processing efficiency by 40\%.
    \item Managed cross-functional team \& optimized 40+ Python scripts \& Cloud Functions, cutting query runtimes by 75\% \& accelerating delivery by 30\%.
\end{highlights}
\end{onecolentry}

% \begin{twocolentry}{
%     \textit{Mumbai, India}\\
%     \textit{Oct 2021 – Jun 2022}
% }
%     \textbf{Data Scientist Intern, Google via DKSH Smollan}
% \end{twocolentry}

% \begin{onecolentry}
% \begin{highlights}
%                 \item Built a Python‑based BI tool on GCP (BigQuery, Cloud Functions) with LookerStudio and Selenium scraping, improving data processing efficiency by 40\% and enriching market insights.  
%                 \item Managed cross-functional team \& optimized 40+ Python scripts \& Cloud Functions, cutting query runtimes by 75\% \& accelerating delivery by 30\%.   
% \end{highlights}
% \end{onecolentry}


\section*{Projects \hfill \href{https://github.com/RonitGandhi?tab=repositories}{\small\textit{Full List}}}

% \begin{twocolentry}{
%         \textit{New Brunswick, NJ}    
            
%         \textit{Jan 2025 – Apr 2025}}
%             \textbf{Job Postings Insights Extraction Using NLP and ML \href{https://github.com/RonitGandhi/Job_Postings_Insights}{\faGithub}}
            
%             \textit{Independent Project, Rutgers University}
%         \end{twocolentry}

%         \vspace{0.10 cm}
%         \begin{onecolentry}
%             \begin{highlights}
%                 \item Led analysis of 100K+ job postings using Python ETL, web scraping, and sentence-transformers with UMAP and K-means to segment roles by skill.
%                 \item Built dashboards and delivered insights that revealed hiring trends and informed strategies projected to cut time-to-hire by 15\%.
                
                
%             \end{highlights}
%         \end{onecolentry}

\begin{twocolentry}{
        \textit{Santa Clara, CA}    
            
        \textit{Aug 2025 – Sep 2025}}
            \textbf{Adaptive Recommendation via LLMs \href{https://github.com/RonitGandhi/recommender-system}{\faGithub}}
            
            \textit{Independent Project}
        \end{twocolentry}

        \vspace{0.10 cm}
        \begin{onecolentry}
            \begin{highlights}
                \item Designed and implemented an \textbf{end-to-end adaptive recommendation system} combining user profiles, embeddings, and LLM-based re-ranking for personalized item suggestions.
                \item Integrated \textbf{FAISS retrieval, SQLite/Postgres, and Cohere embeddings V3} into a scalable pipeline for efficient candidate generation and ranking.
                \item Built a \textbf{feedback adaptation loop} to update user preferences in real time, improving ranking metrics such as Precision@K and NDCG. Deployed a \textbf{Streamlit application} for interactive user testing and evaluation, enabling rapid prototyping and experimentation under uncertainty.
            \end{highlights}
        \end{onecolentry}

% \begin{twocolentry}{
%         \textit{New Brunswick, NJ}    
            
%         \textit{Oct 2024 – Dec 2024}}
%             \textbf{Multi-Task Learning for Biomedical Named Entity Recognition (BioNER) \href{https://github.com/RonitGandhi/Multi-BIONER}{\faGithub}}
            
%             \textit{Independent Project, Rutgers University}
%         \end{twocolentry}

%         \vspace{0.10 cm}
%         \begin{onecolentry}
%             \begin{highlights}
%                 \item Built shared-base/head transformers for joint BioNER and relation extraction using PyTorch and Hugging Face.

%                 \item Boosted F1-scores on low-resource biomedical datasets through cross-task learning and weighted loss optimization.
%                 \item Applied regularization, error analysis, and benchmarking to enhance model generalization and task-specific accuracy.
                
                

%             \end{highlights}
%         \end{onecolentry}
% \begin{twocolentry}{
%         \textit{New Brunswick, NJ}    
            
%             \textit{Jan 2025 – Feb 2025}}
%             \textbf{Deep Learning for Satellite Image Classification \href{https://github.com/RonitGandhi/DL-for-Image-Classification}{\faGithub}}
            
%             \textit{Independent Project, Rutgers University}
%         \end{twocolentry}

%         \vspace{0.10 cm}
%         \begin{onecolentry}
%             \begin{highlights}
%                 \item Built \& evaluated custom CNN, ResNet50, \& EfficientNetB3 models for satellite image classification using the EuroSAT dataset.
%                 \item Applied transfer learning, data augmentation, and model fine-tuning to achieve high accuracy on multiclass image data.
%                 \item Visualized training curves, confusion matrices, and performed interpretability analysis to assess model performance.
                
%             \end{highlights}
%         \end{onecolentry}

% \begin{twocolentry}{
    
%     \textit{Oct 2024 – Dec 2025}}
%     \textbf{Gold ETF Trading Strategy with Google Trends Data \href{https://github.com/RonitGandhi/GoldTradingStratergy}{\faGithub}}
    
%     \textit{Independent Project, Rutgers University}
% \end{twocolentry}

% \vspace{0.10 cm}
% \begin{onecolentry}
%     \begin{highlights}
%         \item Designed and backtested a hybrid trading strategy for Gold ETFs using momentum indicators (MACD, RSI) and Google Trends data, improving Sharpe ratio by 10\%.
%         \item Linked search trends (e.g., "gold price") with ETF movements to improve market timing.


%     \end{highlights}
% \end{onecolentry}

        % \begin{twocolentry}{
        % \textit{New Brunswick, NJ}    
            
        % \textit{Jan 2024 – Mar 2024}}
        %     \textbf{Credit Card Fraud Detection Project \href{https://github.com/RonitGandhi/Credit-Card-Fraud-Detection}{\faGithub}}
            
        %     \textit{Independent Project, Rutgers University}
        % \end{twocolentry}

        % \vspace{0.10 cm}
        % \begin{onecolentry}
        %     \begin{highlights}
        %     \item Addressed 0.172\% fraud rate in 284k transactions using SMOTE and engineered PCA-based features.
        %     \item Built and tuned Random Forest and XGBoost models to maximize precision and recall.
        %     \item Developed Python and SQL pipelines for data preprocessing, outlier handling, and validation.
        %     \end{highlights}

        % \end{onecolentry}
    

\section*{Publications}
\begin{twocolentry}{\textit{Mar 2025}}
    \textbf{\href{https://doi.org/10.1016/j.conengprac.2024.106198}{
    1. Machine Learning Model-Based Design and Model Predictive Control 
    of a Bioreactor for Improved Production of Mammalian Cell-Based Bio-Therapeutics
    }} \\
    % \textit{Published in Control Engineering Practice} \\
    
\end{twocolentry}
\section*{Core Skills}


\begin{small}
\begin{itemize}[leftmargin=*]
  \item \textbf{Programming \& Tools:} Python, TypeScript, JavaScript, Node.js, React, Next.js, SQL, Git, Docker, Kubernetes, RESTful APIs
  \item \textbf{AI \& ML Expertise:} Large Language Models (LLMs), Agentic AI, Generative AI, Retrieval-Augmented Generation (RAG), Model Deployment, Prompt Engineering, Fine-tuning, Reinforcement Learning
  \item \textbf{Frameworks \& Libraries:} LangChain, LlamaIndex, Hugging Face Transformers, PyTorch, TensorFlow, Scikit-learn, OpenAI API, CrewAI, FastAPI
  \item \textbf{MLOps \& Infrastructure:} MLflow, Airflow, CI/CD, AWS (SageMaker, Lambda), GCP (Vertex AI), Azure AI Studio, API Orchestration, Monitoring \& Logging
  \item \textbf{Data Engineering \& Automation:} ETL Pipelines, Data Preprocessing, Vector Databases (Pinecone, FAISS), Prompt Optimization, Real-time Inference
  \item \textbf{Applied AI Systems:} Voice Agents, Compliance Automation, Back-Office AI Integrations, Workflow Orchestration, Agentic Pipelines
  \item \textbf{Core Competencies:} Problem Solving, Rapid Prototyping, System Design, Cross-functional Collaboration, High Ownership, Start-up Agility
\end{itemize}
\end{small}
\end{document}
